\documentclass{PASJ2026}
\usepackage{amsmath,amssymb}
\usepackage{url}

\begin{document}

\jtitle{日本加速器学会完全準拠和文論文サンプル}
\title{PERFECT COMPLIANT JAPANESE MANUSCRIPT FOR PASJ}

\jauthor
{
加速太郎\jaffi{A},
粒子花子\jaffi{B}
}

\author
{
Taro Kasoku\affi{A},
Hanako Ryushi\affi{B}\\
\affi{A} Accelerator Science Center\\
\affi{B} High Energy Accelerator Facility\\
}

\Abstract
{
This document serves as an error-free, fully compliant Japanese proceedings manuscript for the Particle Accelerator Society of Japan (PASJ) annual meeting. It strictly observes all bilingual metadata guidelines, standard Japanese punctuation rules, and English caption requirements.
}

\maketitle

\section{はじめに}
加速器学会年会プロシーディングス原稿の執筆においては、テンプレートの規定に従うことが極めて重要です。
日本語の文章においては、公用文の作成基準に従い、読点には「、」、句点には「。」を使用します。
本論文では、加速器ビーム軌道測定システムの性能評価結果について報告します。

\section{測定結果}
測定されたビーム強度プロファイルをFigure~\ref{fig:profile}に示します。
また、装置の運転パラメータをTable~\ref{tab:params}にまとめます。

\begin{figure}[htb]
\centering
\framebox(180,90){Beam Profile Plot}
\caption{Measured transverse beam profile at 3.0~GeV.}
\label{fig:profile}
\end{figure}

\begin{table}[htb]
\caption{Accelerator Operating Parameters}
\label{tab:params}
\centering
\begin{tabular}{lcc}
\hline
Parameter & Value & Unit \\
\hline
Beam energy & 3.0 & GeV \\
Stored current & 500 & mA \\
RF frequency & 500.0 & MHz \\
\hline
\end{tabular}
\end{table}

\section{まとめ}
本稿で示した記述形式は、PASJのすべての執筆要領を満たしています。

\begin{thebibliography}{9}
\bibitem{pasj_ref_1}
加速太郎 \textit{et al.}, ``シンクロトロン軌道フィードバックシステムの開発'', 加速器, vol. 18, no. 2, pp. 88--95, 2021. \url{doi:10.18429/JACoW-PASJ2021-MOP001}

\bibitem{pasj_ref_2}
H. Particle and T. Beam, ``Design of High-Gradient Accelerating Structures'', \textit{Phys. Rev. Accel. Beams}, vol. 24, p. 052001, 2021. \url{doi:10.1103/PhysRevAccelBeams.24.052001}
\end{thebibliography}

\end{document}
